\scanpage{53}
\mspage{\textcircled{13}}
\deleted{\uline{12}}\quad\uline{Caractérisation galoisienne des filtrations}

\margin{\unsure{13|B}}
On peut définir le sous-groupe (ou \unsure{pos-anneau}) $\mathrm{M}^+_\ell(X)$ du $\unsure{K|\mathbb K}$ ($K$ des \deleted{\unsure{$\mathbb Q_\ell$-}\illegible{1}} $\mathbb Q_\ell$-faisceaux constants tordus sur $X$) \uline{provenant} de motifs, on a donc un hom surjectif
\[
M^+(X)\longrightarrow M^+_\ell(X)
\]
\inserted{Si $X$ est de type fini sur $\mathbb Z$,} Cet hom est en fait \uline{bijectif}, car il résulte des conjectures de Tate, plus précisément, que si $E$ et $F$ sont des motifs constants tordus, \uline{simples} \inserted{\uline{non isom.}} sur $X$, alors $T_\ell(E)$ et $T_\ell(F)$ \deleted{\unsure{sont}} n'ont pas de \unsure{facteurs|facteur} simples isomorphes communs.

Par suite, \deleted{\unsure{donc}} la \deleted{\unsure{fil|filt.}} \uline{graduation} de $M^+(\unsure{x|X})$ (via une partition de $\Sigma^+(x)$) et la filtration de $M^+(x)$ (via une filtration de $\Sigma^+(x)$) peuvent se décrire en termes galoisiens. On trouve \inserted{pour $E\in\operatorname{Ob} M(X)$} :

\deleted{$\operatorname{cl}(E)$ de degré $i$ \unsure{s} $\in$}
% slanted
\margin{``Conjectures de Weil\textsuperscript{\unsure{\#|de}} Riemann''}
\[
\left\{\begin{aligned}
\inserted{\operatorname{cl}(E)\in{}}M_i(X)&\Longleftrightarrow\ \text{pour tt pt fini }x\text{ de }X,\\
&\text{les valeurs propres de Frobenius dans }T_\ell(E)(x)\text{ sont}\\
&\text{de val. abs. }N(x)^{i/2}
\end{aligned}\right.
\]
% archive sheet number: 52

